% DO -- os/samir/samir_main.c repl_try_do_file (DO <name> at the dot prompt
% runs <name>.PRG from the disk; 8191-byte cap, #18) and os/samir/cmd/proc.c
% p_do / p_call (a PROCEDURE of the program running; WITH arguments by value;
% #16 unknown target).  Gates: test-samir-canon-y2k, test-interp-proc.
\begin{XbaseEntry}{DO}{command}
Runs a program file, or a procedure of the program that is running.
\Syntax
\SyntaxLine{\kw{DO} \xp{program}}
\par or\par
\SyntaxLine{\kw{DO} \xp{procedure} \opt{\kw{WITH} \xp{expression list}}}
\Usage
At the dot prompt, \cmd{DO} \xp{program} reads the program file
\xp{program}\texttt{.PRG} from the disk and runs it.  You need not type the
extension.  The program runs from its first line until it reaches
\cmd{RETURN}, its first \cmd{PROCEDURE} or its last line; then the dot prompt
returns.  A program file may hold at most 8,191 characters; a larger file is
refused with \xmsg{18}{Line exceeds maximum of 254 characters.}

In a program, \cmd{DO} \xp{procedure} runs a procedure of the same program
file: the lines that follow its \cmd{PROCEDURE} line, up to the next
\cmd{RETURN}.  When the procedure returns, the program goes on with the line
after the \cmd{DO}.  Procedures may call other procedures, and themselves.  At
most 20 \cmd{DO} commands may be in progress at once, counting the program
started at the dot prompt; the \cmd{DO} that would begin a twenty-first is
refused with \xmsg{103}{DOs nested too deep.}

\kw{WITH} passes the values of the expressions to the procedure, which
receives them in the memory variables of its \cmd{PARAMETERS} line, in order.
The procedure receives a copy of each value: a change it makes to a parameter
does not change the expression that was passed.
\draft{a character constant written directly in the WITH list is not passed
intact by this version (proc.c); pass a character value
through a memory variable or an expression}

A name that is neither a program file on the disk nor a procedure of the
program running is refused with \xmsg{16}{*** Unrecognized command verb.}
\Tips
The programs that come with InitechBase work on the database already in USE:
open it at the dot prompt with \cmd{USE} before you type \cmd{DO}.
\Examples
\file{RAISE.PRG} calls its procedure \texttt{NEWSAL} with two values:
\ProgramListing{RAISE.PRG}
\DotExample{examples/c-do}
\SeeAlso{PARAMETERS, PROCEDURE, RETURN, PUBLIC, PRIVATE}
\end{XbaseEntry}

